\documentclass[10pt,twocolumn]{article}
\usepackage[T1]{fontenc}
\usepackage[margin=0.55in,columnsep=0.20in]{geometry}
\usepackage{amsmath,amssymb,mathtools}
\usepackage[dvipsnames]{xcolor}
\usepackage{booktabs,tabularx}
\usepackage{enumitem}
\usepackage{microtype}
\usepackage[hidelinks]{hyperref}
\usepackage{annotate-equations}
\newcommand{\peMark}[3][NavyBlue]{\eqnmarkbox[#1]{#2}{#3}}
\newcommand{\peAnnotate}[4][]{\annotate[#1]{#2}{#3}{#4}}
\setlength{\parindent}{0pt}
\setlength{\parskip}{1.5pt}
\setlength{\abovedisplayskip}{4pt plus 1pt minus 1pt}
\setlength{\belowdisplayskip}{4pt plus 1pt minus 1pt}
\setlength{\abovedisplayshortskip}{2pt plus 1pt}
\setlength{\belowdisplayshortskip}{3pt plus 1pt minus 1pt}
\setlength{\jot}{2pt}
\setlist{nosep,leftmargin=*,topsep=1pt,partopsep=0pt}
\renewcommand{\section}[1]{\vspace{2pt}\textbf{#1.}\hspace{0.30em}}
\pagestyle{empty}
\frenchspacing
\raggedbottom
\emergencystretch=1em
\begin{document}
{\bfseries Sample Paper Espresso Digest}\hfill{\footnotesize 2026}\par
{\footnotesize A. Researcher and B. Scientist}\par
\vspace{1pt}\hrule\vspace{2pt}

\section{Problem}
Supervised predictors can interpolate a training set while behaving unpredictably on unseen inputs. The central problem is to select parameters that fit observed examples without using unnecessary capacity. The setting assumes independent pairs $(x_i,y_i)$, a differentiable predictor $f_\theta$, and a loss $\ell$ whose scale is comparable across experiments.

\section{Core contribution}
\textbf{Paper claim:} a single regularized objective exposes the tradeoff between empirical fit and parameter complexity, while validation chooses the tradeoff strength. This fixture does not attribute the construction to a real paper; it exercises the digest's information hierarchy, typography, equation callouts, evidence table, and density checks.

\section{Mechanism}
Optimization updates $\theta$ from minibatch gradients. The first term estimates prediction error; the second shrinks weakly supported directions. Larger $\lambda$ increases bias but can reduce variance. Features must be normalized, because an isotropic parameter penalty is otherwise not invariant to feature units.

\textbf{Assumptions.} Samples are representative; train and evaluation losses use the same target semantics; optimization reaches comparable stationary points; and model selection does not inspect the held-out test set.

\vspace{1.1em}
\begin{equation*}
\peMark[NavyBlue]{objective}{\mathcal{L}(\theta)}=
\peMark[OliveGreen]{signal}{\frac{1}{n}\sum_{i=1}^{n}\ell(f_\theta(x_i),y_i)}+
\peMark[WildStrawberry]{constraint}{\lambda\lVert\theta\rVert_2^2}.
\end{equation*}
\peAnnotate[yshift=1.2em]{above}{objective}{objective}
\peAnnotate[yshift=-1.1em]{below}{signal}{data fit}
\peAnnotate[yshift=1.2em]{above}{constraint}{regularizer}
\par\vspace{2.4em}

\textbf{Symbols.} $n$ is sample count; $\theta$ contains trainable parameters; $\ell$ is per-example loss; and $\lambda\geq0$ controls shrinkage. The gradient is the sum of the data gradient and $2\lambda\theta$, so the regularizer continuously pulls parameters toward zero.

\textit{Intuition:} the data term asks for a model that explains observations, whereas the penalty charges for explanations that rely on large coefficients. The optimum is therefore a negotiated point, not the independent optimum of either term.

\section{Selection protocol}
Split data before tuning. For each candidate $\lambda$, fit on the training subset with the same optimizer budget, score on validation data, then refit the chosen configuration without consulting the test labels. Report the complete selection range and random seeds. A useful sanity check is that $\lambda=0$ recovers the unregularized baseline and a very large value approaches a nearly constant predictor.

\section{Evidence}
The synthetic measurements below are compile fixtures, not empirical claims. They encode the comparison shape a digest should preserve: direction, scale, and cost.

\begin{center}
\begin{tabularx}{\columnwidth}{@{}Xrrr@{}}
\toprule
Setting & Train & Test & $\lVert\theta\rVert_2$ \\
\midrule
No penalty & 0.08 & 0.31 & 14.2 \\
Selected $\lambda$ & 0.12 & \textbf{0.19} & 4.8 \\
Strong penalty & 0.28 & 0.30 & 0.9 \\
\bottomrule
\end{tabularx}
\end{center}

The middle row illustrates the intended claim pattern: a modest loss of training fit accompanies improved evaluation error and reduced norm. The strong setting underfits. A serious digest would attach every number to a table, figure, or appendix location in the source paper and would distinguish mean performance from uncertainty.

\section{Failure conditions}
The penalty can be ineffective when scale is unnormalized, the useful solution truly needs large coefficients, distribution shift dominates estimation error, or validation data are too small. Correlated features also make individual parameter magnitudes hard to interpret. In overparameterized networks, equal norms need not imply equal functions, so the regularizer is a training bias rather than a universal complexity measure.

\section{Diagnostics}
Check training and validation curves over $\lambda$; rerun across seeds; compare matched optimization budgets; inspect whether the selected point lies at a search boundary; and report the unregularized baseline. If every setting performs identically, verify that the penalty participates in gradient computation and is not accidentally applied to frozen parameters.

\section{Related concepts}
Ridge regression has a closed-form analogue for squared loss. Under a Gaussian prior, the same penalty appears as maximum a posteriori estimation. Weight decay can match $L_2$ regularization for basic gradient descent but differs for some adaptive optimizers. Early stopping, data augmentation, and architectural constraints impose other forms of inductive bias and should not be treated as interchangeable without evidence.

\section{Reporting checklist}
A faithful summary must state the task, dataset split, model family, loss, regularizer, selection rule, baseline, primary metric, and uncertainty. Record whether the reported result is a mean, a best run, or a single trial. Preserve the direction of every comparison and keep absolute values when they affect practical meaning. If the source omits a detail, say that it is unreported instead of supplying a plausible default.

The digest should also connect claims to evidence. A method claim needs the defining mechanism or equation; a performance claim needs the evaluation setting and comparator; a robustness claim needs the perturbation or distribution tested. Explanatory intuition may translate the paper's notation, but it must remain visibly separate from what the authors actually establish.

\section{Takeaway}
Regularization is useful only relative to a specified data regime, selection protocol, and evaluation measure. The compact objective is the organizing idea; assumptions, quantitative comparisons, and failure modes determine whether the resulting conclusion transfers beyond the reported experiment.

{\footnotesize\textbf{Source:} Local compile and layout-density fixture.}
\end{document}
